% Lean module: MonadFormal.Quotient
\subsubsection{Postulate II: anticipatory compression}\label{ss:anticipatory-compression}\label{ss:relative-compression}

The requirement is to preserve the distinctions admitted continuations
can demand. For a declared family, its exact set-based realization is a
sufficient representation.

A map $f:X\to O$ \emph{factors through} $q:X\to Y$ if there is
$\bar f:Y\to O$ with $f=\bar f\circ q$. A representation $q$ is
sufficient for $F$ if every probe factors through it:
\begin{equation}
  \forall i\in I\;\exists\bar f_i:Y\to O_i,
  \qquad f_i=\bar f_i\circ q.\label{eq:sufficient}
\end{equation}
\leandef{MonadFormal.Factors,MonadFormal.Sufficient}
This is exact preservation of declared distinctions. It is not, by itself,
minimal storage, efficient evaluation or a physically obtainable encoding.

\begin{theorem}[Coarsest sufficient quotient]\label{thm:coarsest-quotient}
Let $Q_F=X/{\approx_F}$ and $q_F(x)=[x]_F$. Then $q_F$ is
surjective and sufficient, and $q_F(x)=q_F(y)$ exactly when
$x\approx_F y$. For every sufficient surjection $q:X\to Y$, there
is a unique $h:Y\to Q_F$ such that $q_F=h\circ q$.
\end{theorem}
\begin{proof}
Define $\bar f_i([x]_F)=f_i(x)$. Observational equivalence makes
this independent of the representative, so each probe descends to
$Q_F$. Every quotient class has a representative. Conversely,
sufficiency implies that $q(x)=q(y)$ forces $f_i(x)=f_i(y)$ for
every $i$, hence $q_F(x)=q_F(y)$. Define $h(q(x))=q_F(x)$.
The preceding implication makes this well-defined, and surjectivity
of $q$ makes it both everywhere defined and unique.
\leanproof{MonadFormal.toQuot_sufficient,MonadFormal.toQuot_eq_iff,MonadFormal.toQuot_surjective,MonadFormal.quotient_universal}
\end{proof}

\begin{theorem}[Unique descent and relabeling]\label{thm:unique-descent}
If $q:X\to Y$ is surjective and $f:X\to O$ is constant on its
fibers, then $f=\bar f\circ q$ for a unique $\bar f:Y\to O$.
If surjections $q:X\to Y$ and $r:X\to Z$ have the same fibers,
their values differ only by a bijection $h:Y\to Z$ satisfying
$r=h\circ q$.
\end{theorem}
\begin{proof}
For each $y$, choose $x$ with $q(x)=y$ and set $\bar f(y)=f(x)$.
Fiber constancy removes dependence on the choice; evaluating any
other candidate at $q(x)$ proves uniqueness. Apply this construction
to $r$ through $q$, and to $q$ through $r$. The two resulting maps
compose to the identity on the images, which are the whole codomains.
\leanproof{MonadFormal.unique_descent,MonadFormal.same_fibers_relabeling}
\end{proof}

No finiteness hypothesis is needed. The construction can be
noncomputable: existence of a quotient is not an algorithm for
recognizing its classes.
